\documentclass[11pt,a4paper]{article}
\usepackage[margin=1in]{geometry}
\usepackage{amsmath,amssymb,amsthm}
\usepackage{booktabs}
\usepackage{enumitem}
\usepackage[most]{tcolorbox}
\usepackage{hyperref}

\newtcolorbox{keybox}[1][]{colback=blue!4,colframe=blue!50!black,fonttitle=\bfseries,title=#1}
\newtcolorbox{pitfall}[1][]{colback=red!4,colframe=red!60!black,fonttitle=\bfseries,title=#1}
\newtcolorbox{interview}[1][]{colback=green!4,colframe=green!40!black,fonttitle=\bfseries,title=#1}

\DeclareMathOperator*{\argmax}{arg\,max}
\DeclareMathOperator*{\argmin}{arg\,min}
\newcommand{\E}{\mathbb{E}}
\newcommand{\Var}{\mathrm{Var}}

\title{MLE and MAP: Interview Prep Notes}
\author{Notes from a back-and-forth lecture}
\date{\today}

\begin{document}
\maketitle
\tableofcontents
\newpage

%--------------------------------------------------------------
\section{Prerequisites (quick mental check)}
%--------------------------------------------------------------
\begin{itemize}[nosep]
  \item \textbf{Probability basics:} random variables, PDFs/PMFs, independence. The i.i.d.\ assumption lets the likelihood factor into a product.
  \item \textbf{Common distributions:} Bernoulli, Binomial, Gaussian, (Poisson).
  \item \textbf{Logs and calculus:} product $\to$ sum, derivative set to zero.
  \item \textbf{Likelihood vs probability:} $p(\text{data}\mid\theta)$ viewed as a function of $\theta$ with data fixed; \emph{not} a distribution over $\theta$.
  \item \textbf{Expectation and variance:} needed for the bias discussion.
\end{itemize}
Optional later: KL divergence (MLE $\equiv$ minimizing KL between empirical distribution and model), Fisher information (asymptotic normality, Cram\'er--Rao).

%--------------------------------------------------------------
\section{Notation}
%--------------------------------------------------------------
\begin{center}
\begin{tabular}{@{}ll@{}}
\toprule
Symbol & Meaning \\
\midrule
$x_1,\dots,x_n$ & the $n$ observed data points \\
$\theta$ & unknown model parameter \\
$p(x\mid\theta)$ & probability (or density) of $x$ when the parameter is $\theta$ \\
$L(\theta)$ & likelihood, $L(\theta)=p(x_1,\dots,x_n\mid\theta)$, as a function of $\theta$ \\
$\prod$ & product, $\prod_{i=1}^n a_i = a_1 a_2\cdots a_n$ \\
$\argmax_\theta f(\theta)$ & the \emph{value} of $\theta$ that maximizes $f$ \\
$\hat\theta$ & an estimate of $\theta$ computed from data \\
\bottomrule
\end{tabular}
\end{center}

\paragraph{i.i.d.} \emph{Independent}: one point tells nothing about another, so the joint factors into a product ($P(A\cap B)=P(A)P(B)$). \emph{Identically distributed}: all points share the same distribution and the same $\theta$.

%--------------------------------------------------------------
\section{MLE: the core idea}
%--------------------------------------------------------------
MLE asks: \emph{which parameter value makes the data I actually saw least surprising?}
For i.i.d.\ data,
\[
L(\theta)=\prod_{i=1}^n p(x_i\mid\theta),\qquad
\hat\theta_{\text{MLE}}=\argmax_\theta L(\theta).
\]

\begin{keybox}[Likelihood vs probability]
$L(\theta)$ has the same formula as the probability of the data, but $\theta$ varies and the data is fixed. It does \textbf{not} integrate to 1 over $\theta$ and is \textbf{not} ``the probability that $\theta=0.7$''. That statement needs a prior (MAP/Bayes).
\end{keybox}

\subsection{Why maximize the log-likelihood}
Since $\log$ is strictly increasing, $\argmax_\theta \log L(\theta)=\argmax_\theta L(\theta)$ (the maximizer keeps the same location in $\theta$). Reasons to use it:
\begin{itemize}[nosep]
  \item \textbf{Product $\to$ sum:} $\log\prod_i p(x_i\mid\theta)=\sum_i\log p(x_i\mid\theta)$, easier to differentiate.
  \item \textbf{Numerical stability:} a product of many numbers $<1$ underflows to $0$; a sum of logs does not.
\end{itemize}

\begin{pitfall}[Common mistake]
Taking the log does \emph{not} create uniqueness of the maximizer; $L$ and $\log L$ have identical maximizers. (Log-concavity of a model can help, but it is a separate property.)
\end{pitfall}

\subsection{Recipe}
\begin{enumerate}[nosep]
  \item Write $L(\theta)=\prod_i p(x_i\mid\theta)$.
  \item Take the log; drop any terms without $\theta$.
  \item Differentiate, set to $0$, solve.
  \item Check the second derivative is negative (a max, not a min), and that the answer lies in the valid parameter range.
\end{enumerate}

\begin{pitfall}[Sanity check]
If $\theta$ is a probability and you get $\hat\theta=2$, there is an algebra slip. Re-check the chain rule: $\frac{d}{d\theta}\log(1-\theta)=-\frac{1}{1-\theta}$.
\end{pitfall}

%--------------------------------------------------------------
\section{MLE for the Bernoulli (coin)}
%--------------------------------------------------------------
Sequence H,H,T with $\theta=P(\text{heads})$: $L(\theta)=\theta^2(1-\theta)$, so $\log L=2\log\theta+\log(1-\theta)$ and
\[
\frac{2}{\theta}-\frac{1}{1-\theta}=0\;\Rightarrow\;2(1-\theta)=\theta\;\Rightarrow\;\hat\theta=\tfrac23 .
\]
Second derivative $-\frac{2}{\theta^2}-\frac{1}{(1-\theta)^2}<0$, so $\log L$ is concave and $2/3$ is the unique maximum.

\paragraph{General case: $k$ heads in $n$ flips.}
\[
\log L(\theta)=k\log\theta+(n-k)\log(1-\theta)
\;\Rightarrow\;
\frac{k}{\theta}=\frac{n-k}{1-\theta}
\;\Rightarrow\;
\boxed{\hat\theta_{\text{MLE}}=\frac{k}{n}}
\]

\paragraph{Where is $\binom{n}{k}$?} If you record only the count $k$, $P(k\mid\theta)=\binom nk\theta^k(1-\theta)^{n-k}$ and
$\log L=\log\binom nk+k\log\theta+(n-k)\log(1-\theta)$. The term $\log\binom nk$ does not contain $\theta$, so it vanishes on differentiation. This is exact (no Stirling approximation). Likewise, for a specific sequence there is no coefficient at all. \textbf{General rule:} any factor of $L(\theta)$ independent of $\theta$ can be dropped.

%--------------------------------------------------------------
\section{MLE for the Gaussian}
%--------------------------------------------------------------
For i.i.d.\ $x_i\sim\mathcal N(\mu,\sigma^2)$,
\[
p(x\mid\mu,\sigma^2)=\frac{1}{\sqrt{2\pi\sigma^2}}\exp\!\Big(-\frac{(x-\mu)^2}{2\sigma^2}\Big),
\]
\[
\log L(\mu,\sigma^2)=-\frac n2\log(2\pi\sigma^2)-\frac{1}{2\sigma^2}\sum_{i=1}^n(x_i-\mu)^2 .
\]

\paragraph{Mean.} $\mu$ appears only in the last term, so maximizing over $\mu$ means minimizing the \textbf{sum of squared errors} $\sum_i(x_i-\mu)^2$. Differentiating: $-2\sum_i(x_i-\mu)=0$, so
\[
\hat\mu_{\text{MLE}}=\frac1n\sum_i x_i\quad\text{(the sample mean).}
\]
(Wording: it is the mean of the \emph{sample}, an estimate of the unknown true $\mu$.)

\paragraph{Variance.} With $v=\sigma^2$ and $\mu=\hat\mu$:
\[
-\frac{n}{2v}+\frac{1}{2v^2}\sum_i(x_i-\hat\mu)^2=0
\;\Rightarrow\;
\boxed{\hat\sigma^2_{\text{MLE}}=\frac1n\sum_i(x_i-\hat\mu)^2}
\]

\subsection{Bias: why the sample variance uses $n-1$}
\begin{keybox}[Intuition (lead with this in interviews)]
You used the data to compute $\bar x$, so the deviations are forced to satisfy $\sum_i(x_i-\bar x)=0$. Only $n-1$ of them are free: $n-1$ degrees of freedom. Also, $\bar x$ minimizes $\sum_i(x_i-c)^2$ over $c$, so
$\sum_i(x_i-\bar x)^2\le\sum_i(x_i-\mu)^2$, which makes dividing by $n$ underestimate $\sigma^2$.
\end{keybox}

\paragraph{Exact calculation.} Using $\sum_i(x_i-\mu)^2=\sum_i(x_i-\bar x)^2+n(\bar x-\mu)^2$:
\begin{align*}
\E\Big[\sum_i(x_i-\mu)^2\Big]&=n\sigma^2,\qquad
\E\big[n(\bar x-\mu)^2\big]=n\Var(\bar x)=\sigma^2,\\
\Rightarrow\;\E\Big[\sum_i(x_i-\bar x)^2\Big]&=(n-1)\sigma^2 .
\end{align*}
Hence dividing by $n-1$ is unbiased, while $\E[\hat\sigma^2_{\text{MLE}}]=\frac{n-1}{n}\sigma^2$.

\begin{itemize}[nosep]
  \item The MLE is biased but \textbf{consistent}: the bias $\to0$ as $n\to\infty$.
  \item The $n-1$ estimate is the larger of the two. For $n=100$ the ratio is $100/99\approx1\%$ (small); for $n=5$ it is $25\%$.
  \item The sample \emph{standard deviation} is still slightly biased (concavity of $\sqrt{\cdot}$, Jensen).
\end{itemize}

%--------------------------------------------------------------
\section{MLE and loss functions}
%--------------------------------------------------------------
\begin{keybox}[Core principle]
\textbf{Minimizing a loss = maximizing a log-likelihood under some noise assumption.}
Loss $=$ negative log-likelihood (NLL), purely because optimizers are conventionally written to \emph{minimize}.
\end{keybox}

\subsection{Gaussian noise $\Rightarrow$ MSE}
For $y_i=f(x_i;\theta)+\varepsilon_i$, $\varepsilon_i\sim\mathcal N(0,\sigma^2)$:
\[
\log L(\theta)=\text{const}-\frac{1}{2\sigma^2}\sum_i\big(y_i-f(x_i;\theta)\big)^2,
\]
so maximizing $\log L$ $\Leftrightarrow$ minimizing the squared error (MSE up to the factor $1/n$).

\subsection{Laplace noise $\Rightarrow$ MAE}
Density $p(y\mid f)=\frac{1}{2b}\exp(-|y-f|/b)$:
\[
\log L(\theta)=-n\log(2b)-\frac1b\sum_i|y_i-f(x_i;\theta)| ,
\]
so MLE minimizes $\sum_i|y_i-f(x_i;\theta)|$, the absolute-error (L1) loss.

\subsection{Location estimate: mean vs median}
\begin{center}
\begin{tabular}{@{}lll@{}}
\toprule
Noise & Loss & Location estimate \\
\midrule
Gaussian & squared & mean \\
Laplace & absolute & median \\
\bottomrule
\end{tabular}
\end{center}
\emph{Why median:} nudging $m$ right increases the distance to points on the left and decreases it to points on the right; $\sum_i|x_i-m|$ stops improving when half the points lie on each side. Example: $\{1,2,3,4,100\}$ has mean $22$ but median $3$. The Laplace density has heavier tails than the Gaussian, so a far point is less surprising and pulls the fit less (linear vs quadratic penalty): MAE is robust to outliers.

\subsection{Bernoulli $\Rightarrow$ binary cross-entropy}
Let $y_i\in\{0,1\}$ and $p_i=f(x_i;\theta)=P(y_i=1\mid x_i)$.
\begin{enumerate}[nosep]
  \item One point (exponents act as a switch):
  $P(y_i\mid x_i;\theta)=p_i^{\,y_i}(1-p_i)^{1-y_i}$. Check: $y_i=1\to p_i$, $y_i=0\to1-p_i$.
  \item Dataset (independence): $L(\theta)=\prod_i p_i^{\,y_i}(1-p_i)^{1-y_i}$.
  \item Log: $\log L=\sum_i\big[y_i\log p_i+(1-y_i)\log(1-p_i)\big]$.
  \item Negate and average (dividing by $n$ does not change the argmin):
\end{enumerate}
\[
\boxed{\mathrm{BCE}(\theta)=-\frac1n\sum_i\big[y_i\log p_i+(1-y_i)\log(1-p_i)\big]}
\]
With a constant $p_i=\theta$ this reduces to the coin formula $k\log\theta+(n-k)\log(1-\theta)$.

\subsection{Categorical $\Rightarrow$ multi-class cross-entropy}
Labels are one-hot $y_{ik}\in\{0,1\}$, $\sum_k y_{ik}=1$; the model gives $p_{ik}=P(\text{class }k\mid x_i)$ (typically via softmax), $\sum_k p_{ik}=1$.
\[
P(y_i\mid x_i;\theta)=\prod_{k=1}^K p_{ik}^{\,y_{ik}},\quad
L(\theta)=\prod_i\prod_k p_{ik}^{\,y_{ik}},\quad
\log L=\sum_i\sum_k y_{ik}\log p_{ik},
\]
\[
\boxed{\mathrm{CE}(\theta)=-\frac1n\sum_i\sum_{k=1}^K y_{ik}\log p_{ik}}
\]
With one-hot labels the inner sum collapses to $-\log p_{i,c}$ ($c$ = true class). Binary cross-entropy is the $K=2$ special case. Examples: assigned probability $0.25\Rightarrow$ loss $-\log0.25=\log4\approx1.386$; probability $1\Rightarrow$ loss $0$; probability $\to0\Rightarrow$ loss $\to\infty$ (confident wrong answers are punished very harshly).

\begin{interview}[Interview: why cross-entropy and not MSE for logistic regression?]
\begin{itemize}[nosep]
  \item BCE is exactly the NLL of a Bernoulli model, which matches a $0/1$ label; MSE assumes Gaussian noise, which does not.
  \item With a sigmoid output, the MSE gradient vanishes when the model is confidently wrong (saturation), so learning stalls. For softmax/sigmoid + CE, the gradient with respect to the logits is simply $p-y$, which stays large exactly when the model is wrong.
\end{itemize}
\end{interview}

%--------------------------------------------------------------
\section{MLE summary and weakness}
%--------------------------------------------------------------
\begin{enumerate}[nosep]
  \item $L(\theta)=p(\text{data}\mid\theta)$ as a function of $\theta$; a product for i.i.d.\ data.
  \item Take the log (product $\to$ sum, stable), maximize, drop terms without $\theta$.
  \item Bernoulli: $\hat\theta=k/n$. Gaussian: sample mean; $\hat\sigma^2$ with $1/n$ (biased, consistent).
  \item Loss = NLL: Gaussian $\to$ MSE, Laplace $\to$ MAE, Bernoulli $\to$ BCE, categorical $\to$ CE.
  \item Properties: consistent, asymptotically normal, invariant under reparametrization.
\end{enumerate}

\paragraph{Weakness: overfitting to small data.} Three flips, three heads: $L=\theta^3$, $\log L=3\log\theta$, derivative $3/\theta$ has no zero. Since $\theta\in[0,1]$, the maximum is at the boundary, $\hat\theta_{\text{MLE}}=1$ (agrees with $k/n$). This says tails are impossible, though a fair coin gives 3 heads with probability $1/8$. MLE trusts the data completely and cannot express prior knowledge. \textbf{MAP fixes this.}

%--------------------------------------------------------------
\section{MAP: the core idea}
%--------------------------------------------------------------
Bayes' rule:
\[
\underbrace{p(\theta\mid D)}_{\text{posterior}}=\frac{\overbrace{p(D\mid\theta)}^{\text{likelihood}}\;\overbrace{p(\theta)}^{\text{prior}}}{p(D)} ,
\]
where $p(D)$ does not depend on $\theta$. Then
\[
\hat\theta_{\text{MAP}}=\argmax_\theta p(\theta\mid D)=\argmax_\theta\big[\log p(D\mid\theta)+\log p(\theta)\big].
\]
So \textbf{MAP = MLE objective + log-prior}.

\begin{itemize}[nosep]
  \item Flat prior $\Rightarrow$ $\log p(\theta)$ is constant $\Rightarrow$ MAP $=$ MLE.
  \item \textbf{Framing:} MLE (frequentist): $\theta$ is a fixed unknown constant, the data is random. MAP (Bayesian): $\theta$ is a random variable with a distribution (prior $\to$ posterior), the data is observed.
  \item MAP is only ``half Bayesian'': it uses a prior but reports a single point (the posterior peak) instead of the full distribution.
\end{itemize}

%--------------------------------------------------------------
\section{MAP for the coin with a Beta prior}
%--------------------------------------------------------------
Prior: $p(\theta)\propto\theta^{\alpha-1}(1-\theta)^{\beta-1}$, with hyperparameters $\alpha,\beta>0$. ($\alpha=\beta=1$ is flat; $\alpha=\beta=2$ is a mild belief that the coin is roughly fair.)

Log-posterior (constants dropped):
\[
J(\theta)=(k+\alpha-1)\log\theta+(n-k+\beta-1)\log(1-\theta).
\]
Let $A=k+\alpha-1$, $B=n-k+\beta-1$, so $J=A\log\theta+B\log(1-\theta)$.
\begin{align*}
\frac{dJ}{d\theta}&=\frac A\theta-\frac B{1-\theta}=0
\;\Rightarrow\;A(1-\theta)=B\theta
\;\Rightarrow\;\theta=\frac{A}{A+B},
\end{align*}
and $A+B=n+\alpha+\beta-2$. Hence
\[
\boxed{\hat\theta_{\text{MAP}}=\frac{k+\alpha-1}{n+\alpha+\beta-2}}
\]
Second derivative $-\frac{A}{\theta^2}-\frac{B}{(1-\theta)^2}<0$ when $A,B>0$, so it is a maximum.

\begin{itemize}[nosep]
  \item \textbf{Check:} $\alpha=\beta=1$ gives $k/n$ (MLE).
  \item \textbf{Example:} $k=n=3$, $\alpha=\beta=2$: $\hat\theta_{\text{MAP}}=\frac{3+2-1}{3+2+2-2}=\frac45=0.8$, versus $\hat\theta_{\text{MLE}}=1$.
  \item \textbf{Pseudo-counts:} the prior acts like $(\alpha-1)$ imaginary heads and $(\beta-1)$ imaginary tails added before the real data. With $\alpha=\beta=2$: $(3+1)/(3+2)=4/5$.
  \item \textbf{Washing out:} as $n\to\infty$, $k$ and $n$ dominate $\alpha,\beta$, so $\hat\theta_{\text{MAP}}\to k/n=\hat\theta_{\text{MLE}}$.
\end{itemize}

%--------------------------------------------------------------
\section{Still to cover}
%--------------------------------------------------------------
\begin{enumerate}[nosep]
  \item \textbf{Regularization link:} Gaussian prior on weights $\to$ L2 (ridge); Laplace prior $\to$ L1 (lasso), derived like MSE/MAE from MLE.
  \item \textbf{MAP vs MLE vs full Bayesian:} when MAP helps (small data, overfitting, high dimension); limits (point estimate; not invariant under reparametrization, unlike MLE).
  \item \textbf{Rapid-fire interview questions} on MLE and MAP together.
  \item Open question: with $\alpha=\beta=2$ and 30 heads in 30 flips, what is $\hat\theta_{\text{MAP}}$? (Use the boxed formula; compare with the 3-flip case.)
\end{enumerate}

\end{document}
